\documentclass[11pt]{article}
\usepackage[margin=24mm]{geometry}
\usepackage{amsmath,amssymb,amsthm,mathtools,hyperref}
\newtheorem{theorem}{Theorem}
\newtheorem{lemma}[theorem]{Lemma}
\newcommand{\R}{\mathbb R}
\newcommand{\eps}{\varepsilon}
\allowdisplaybreaks
\title{An evaluated testing dictionary for the positive-angular C1 main Gram}
\author{}\date{23 October 2026}
\begin{document}\maketitle

This calculation concerns the main Gram of the all-$H$, bounded-$z$
C1 certificate at public commit
\texttt{0c1ce2e15ed067a6845464e2468559cea5274f2c}.
It gives fixed rational bounds on a genuine source/testing rectangle,
including source points approaching $D=0$, $\beta=0$, and $H=+\infty$.
The condition $\lambda>0$ is retained. The separate $\lambda\le0$
construction and the retained squares outside this Gram are not
compiled here. Thus this is not yet a local whole-parent pencil or a
full-set pencil.

\section{The rectangle and source formulas}

Fix the rational constants
\[
 \eps=2^{-160},\qquad \delta=2^{-100},\qquad
 c_0=2^{-8},\qquad A_*=2^{30},\qquad
 \gamma={c_0\over14A_*}={1\over14\,2^{38}}.
\]
The source domain consists of
\[
 0<x,\sigma\le\eps,\quad 0<z\le1,\quad 0<b<1,
 \quad r=bz<1,\quad 0<h<1,
\]
with the original strict gaps and the two nonnegative event
coefficients specified below. Set
\[
 Q=1-b,\quad \theta=1-h,\quad H=h/\sigma,\quad e=xh/\sigma,
 \quad E=1-e>0,\quad k=\sqrt x,\quad q=\sqrt z,
\]
\[
 K=1+x,\quad R_s=1+xr,\quad M=R_s-xz\theta>0,
 \quad J=1+xh-\sigma(1-r)>0,
\]
\[
 P=\sigma+K+\sigma z(1+bx+2xh)
                  +\sigma bx^2hz^2+x^2h\theta z,
\]
\[
\begin{gathered}
 p={z(1+xh)E+(1-r)(1+xhz)\over2P},\qquad
 l={(1+\sigma r)\{eR_s(1+\sigma z)+ME\}\over2P},\\
 \kappa={K(1+\sigma r)(1+\sigma z)\over2MP}.
\end{gathered}
\]
The original coefficients are
\[
\begin{gathered}
 m=\kappa x^2zH=\kappa xze,\qquad
 C={Kz(1+\sigma r)E\over2(1-r)P},\\
 A=p/r,\quad \alpha=l/(zQ),\quad a=p/b,\quad \ell=l/Q,\\
 D={K(1+\sigma z)(1+\sigma r)-JP\over2JP}\ge0,\qquad
 \beta={K(1+\sigma r)(1+xhz)-\theta P\over2\theta P}\ge0.
\end{gathered}
\]
Here $H$ has no upper bound. Neither $b$, $Q$, $\theta$, $1-r$,
nor $z$ is bounded below. The testing rectangle is independently
\[
                  0\le u,w\le\delta.
\]
For signed angles $u_s,w_s$, set
$\alpha_s=\operatorname{sgn}(u_s)$ and
$\beta_s=\operatorname{sgn}(w_s)$, choosing the value $1$ at zero.
The exact testing gauge is
\[
 (u,w,X,Y,Z,W)=(|u_s|,|w_s|,\beta_sX_s,\alpha_sY_s,
                                \alpha_s\beta_sZ_s,W_s).
\]
It preserves the parent polynomial. Composing every testing function
below with this one fixed piecewise semialgebraic map extends the
same feature count to $|u_s|,|w_s|\le\delta$; no square root of a
negative angle is taken. State coordinates $X,Y,Z,W$ are unrestricted
real numbers; no state ratio is bounded.

Define
\[
 H_0=1+zQ,\quad B_s=P+xz(J-\sigma),
\]
\[
 \mathcal H=R_sH_0+z\{(1-r)xh(2+xr)
 -e[2+x(1+h+r\theta)]+Ke^2\}.
\]
Two exact source identities, obtained by multiplication, are
\begin{align}
 2\theta P\beta&=hB_s-\sigma H_0,\label{eq:BE3}\\
 2P(JD+x\theta\beta)&=\sigma^2\mathcal H.\label{eq:BE4}
\end{align}
The angular quantities used below are
\[
\begin{gathered}
 \nu=(2D-2\beta\theta+h)/\sigma,\quad
 \lambda=(8Dh+4\beta\theta^2+2h^2)/\sigma-4bp\sigma z>0,\\
 K_A=4\kappa H-\lambda,\quad d_G=4D/\sigma^2,
 \quad \rho=\lambda x/\sigma,\quad S=d_G+\rho,\\
 L_0=\lambda\sigma,\quad q_0=\lambda k,\quad
 c=\max(0,1-K_A),\quad a_0=k/\sigma.
\end{gathered}
\]
The letter $q_0$ here is a source quantity, not the constant
coefficient of a later scalar channel.

The complete reciprocal parent being certified is
\[
\begin{split}
 F_w={}&4m(W-wX+uY-uwZ)^2+4C(1+w^2)(uW-Y)^2\\
 &+4A(1+u^2)(1+w^2)Y^2
 +(4D+2)(uwW+uX-wY-Z)^2\\
 &+(4\alpha+2)(1+u^2)(Y-wZ)^2
 +(4\beta+2)(1+w^2)(X+uZ)^2\\
 &-2(1+u^2)(1+w^2)Z^2.
\end{split}
\]
All four state coordinates occur; they need not be normalized.

\section{Uniform source bounds, without an upper bound on $H$}

All inequalities in this section follow under the larger condition
$\eps\le2^{-22}$, so in particular on the chosen rectangle.
Directly from the positive formula for $P$ and the definitions,
\[
 1\le P\le1+5\eps,\quad |J-1|,|M-1|\le\eps,
 \quad 1\le B_s\le1+7\eps,\quad 1\le H_0\le2.
\]
Initially $0<e<1$ gives $0\le\mathcal H\le4$: nonnegativity
uses \eqref{eq:BE4}. Eliminate $\beta$ to obtain
\[
 eB_s=xH_0+\sigma\mathcal H-2PJD/\sigma.
\]
Consequently
\begin{equation}\label{eq:basic-bounds}
 e\le2x+4\sigma\le6\eps,\qquad
 |\mathcal H-H_0|\le32\eps,\qquad
 0\le D\le4\sigma^2,\quad 0\le\beta\theta\le1.
\end{equation}
For example, in $\mathcal H-H_0$ the terms are bounded by
$2\eps+3\eps+3e+2e^2$, which is at most $32\eps$ here.
The same positive formulas give
\begin{equation}\label{eq:p-bounds}
 |p-H_0/2|\le16\eps,\quad |l-1/2|\le8\eps,
 \quad 1/4\le p,l,\kappa\le2,
 \quad \kappa\le1,\quad m\le6\eps^2,
 \quad C(1-r)\le1.
\end{equation}
In particular $4a,4\ell\ge1$. The unbounded positive coefficients
$a,\ell,C$ will be retained in whole squares and will not be
estimated separately.

Exact rearrangement gives
\[
 \nu=2D/\sigma+[H_0-ez(J-\sigma)]/P,
\]
\[
 \lambda=2H+2e\theta z(J-\sigma)/P
                  -2\theta H_0/P+8Dh/\sigma-4bp\sigma z.
\]
Thus
\begin{equation}\label{eq:lambda-bounds}
 |\nu|\le3,\qquad |\lambda-2H|\le5,
 \qquad |L_0-2h|\le8\sigma.
\end{equation}

For completeness, the cancellation needed when $H$ is large is
polynomial. Write
\[
 N_\kappa=K(1+\sigma r)(1+\sigma z)-MP=xF_x+\sigma F_\sigma.
\]
Here $F_\sigma=r-1+\sigma rz$, and $F_x$ is the polynomial
$(N_\kappa-N_\kappa|_{x=0})/x$. Its coefficient $\ell^1$ norm
in $(x,\sigma,z,b,h,r)$ is $37$, as can also be checked by
expanding the displayed product. Since all six arguments lie in
$[0,1]$,
\[
 H(2\kappa-1)={eF_x+hF_\sigma\over MP},\qquad
 |H(2\kappa-1)-h(r-1)|\le462\eps.
\]
The estimate $|MP-1|\le8\eps$ suffices for the second inequality.
Combining this with \eqref{eq:lambda-bounds} more precisely yields
\begin{equation}\label{eq:KA-bounds}
 |K_A-\{2\theta H_0-2h(1-r)\}|\le2^{12}\eps,
 \qquad |K_A|\le5,\quad 0\le c\le6.
\end{equation}
If $c>0$, then $h\ge1/8$: for $h<1/8$ the expression in braces
is at least $3/2$, whereas the error is less than $1/4$.
The earlier elimination identity also gives
\[
 a_0^2(hB_s-\sigma H_0)\le\mathcal H.
\]
Therefore, whenever $c>0$,
\begin{equation}\label{eq:c-positive}
 a_0\le8,\qquad L_0\ge1/8,\qquad
 1/\lambda\le8\sigma.
\end{equation}
Every term below containing $a_0$ or $1/\lambda$ from the angular
subtraction is multiplied by $c$. On $c=0$ that subtraction is
identically zero; no uniform estimate on those unused factors is
required.

Put
\[
 R_0=-4\{2(C+D)(r-1)+2p+r-1\}.
\]
The positive expression for $C(1-r)$ and the formula for $p$ give
$|R_0|\le256(x+\sigma)\le512\eps$. Since
\[
 (x+\sigma)^2H=xe+2\sigma e+\sigma h\le2\eps,
\]
\eqref{eq:lambda-bounds} further gives
\begin{equation}\label{eq:R0-bounds}
 |R_0|\sqrt\lambda\le1024\sqrt\eps.
\end{equation}
Finally the exact identity
\[
 S={2\mathcal H\over P}+d_G(1-J+2hx)
       +{2eH_0\over P}-{2e^2z(J-\sigma)\over P}-4bpxz
\]
shows
\begin{equation}\label{eq:S-bounds}
 |S-2H_0|\le256\eps,\qquad
 1\le S\le5,\quad 0\le d_G,\rho\le5,
 \quad 0<L_0\le3,\quad 0<q_0\le4.
\end{equation}
Here $q_0^2=L_0\rho$; all bounds use $\lambda>0$ but no lower
bound on $\lambda$ or upper bound on $H$.

\section{A finite formula for the main Gram}

The state roots are
\[
\begin{gathered}
 N_x=X+uhZ,\quad N_y=(Y-bwZ)/q,\quad
 R=wZ/q-uqW,\quad \Gamma=\sigma Z-uwW,\\
 A_\theta=q(u^2-x)W,\quad T_\theta=kq(u-k)W,
 \quad J_C=N_y+bR,\quad J_\ell=N_y-QR,\quad
 \Gamma_*=\sqrt{d_G}\Gamma.
\end{gathered}
\]
Let
\[
 B_1=\sqrt{\lambda\sigma}\,uZ,\quad
 B_2=\sqrt\rho\,uwW,\quad
 B_3=\sqrt\lambda\,xqW,\quad
 B_4=\sqrt{\lambda x}\,wZ/q.
\]
When $WZ>0$, put
$G_*=\operatorname{sgn}(W)(B_1B_2B_3B_4)^{1/4}$ and
$d_i=B_i-G_*$. This is the positive-state convention followed,
when necessary, by changing the signs of all roots together.
When $WZ\le0$
put $d_i=B_i$. Equivalently, for every state,
\[
 G_*=\sqrt{\lambda x}\,\operatorname{sgn}(W)
                \sqrt{uw\max(WZ,0)},\qquad d_i=B_i-G_*.
\]
Define
\[
 U_*=\sqrt{q_0}\sqrt u\,q(u-k)W,
 \qquad K_*=\sqrt{q_0}\sqrt k\,q(u-k)W.
\]
Use the twelve-root column
\[
 g=(N_x,J_C,J_\ell,\Gamma_*,A_\theta,T_\theta,
                         d_1,d_2,d_3,d_4,U_*,K_*)^{\mathsf T}.
\]
The two radial changes are source-only and uniformly bounded:
\[
 R=J_C-J_\ell,\qquad N_y=QJ_C+bJ_\ell.
\]
Set
\[
 D_*={\rm diag}(2,1,1,1,1/2,1/2,1/2,1/2,1/2,1/2,1,1/2).
\]
We now define a polynomial quadratic form $\mathcal E(g)$, which
fully specifies the correction matrix; a cross term is divided by
two when placed in its two symmetric entries.

For its coefficients put $t=u^2$, $\xi=w^2$, $g_s=h-Q$ and
\[
\begin{gathered}
 R_x=-4\{2Db+(1+2m)g_s-2\beta\theta\},\quad B_x=8\kappa e,\\
 A_y=8(p-l)+4z\{(1+2m)g_s+2Dh+2p\sigma\},\\
 \chi=8(pb+lQ)-4h+4eh(2D+1)-8Dh\\
 \hspace{12mm}+4z\{4Dbh+4pb\sigma+(1+2m)g_s^2+2\beta\theta^2\},\\
 k_1=2(1+2m)g_s^2+8Dhb+8pb\sigma+4\beta\theta^2.
\end{gathered}
\]
Then
\[
\begin{split}
 E_{yy}&=2z\{(1+2m)u^2+(1+2D)w^2\},\\
 E_{RR}&=z\{8bD+8bp\sigma+4b\sigma(2C+2D+1)(r-1)
                                  +k_1u^2+2b^2(2D+1)w^2\},\\
 E_{yR}&=4z\{u^2[2Dh+2p\sigma+(1+2m)g_s]+b(2D+1)w^2\\
 &\hspace{25mm}+\sigma(2C+1)(r-1)+2D[1+\sigma(r-1)]+2p\sigma\}.
\end{split}
\]
The entire first normal is represented by the linear form
\[
 F_1={\sqrt{d_G}\over S}u\Gamma_*
 +{\sqrt\rho\over S}\{kd_1-ud_2+(u-k)d_3\}
                         -{\rho\over S}\sqrt\sigma T_\theta.
\]
The entire second normal is represented by
\[
\begin{split}
 F_2={}&{R_0\sqrt{d_G}\over S}qw\Gamma_*
 +{R_0\rho\sigma z\over S}R
 -{R_0\sqrt\rho\over S}(qw+z\sqrt\sigma)(d_2-d_3)\\
 &+\left\{{R_0q_0z\over S}
       +{R_0\sqrt\lambda\sqrt\rho\over S}qw\right\}T_\theta
                                          -8pquwF_1.
\end{split}
\]
These are exact identities $F_1=u\Gamma$ and
$F_2=qw(R_0-8pu^2)\Gamma$. In particular they keep both complete
state-normal couplings, not only their contact restrictions.

Write
\[
 \eps_C=2/\lambda+u+k,\qquad
 \eps_G=2a_0u+(u+k)/L_0.
\]
The correction is exactly
\begin{align*}
 \mathcal E={}&E_{yy}N_y^2+E_{RR}R^2+E_{yR}N_yR
                  -4\nu N_xF_1+J_CF_2\\
 &-4uqw(2m+2D+1)N_xN_y+uqwR_xN_xR
                 +qwB_xN_xA_\theta+uA_yN_yA_\theta+\chi uA_\theta R\\
 &-\tfrac12\rho\sigma R^2+\sqrt{q_0}\sqrt u\,U_*R
       -{cd_G\over S}\eps_C\sum_{i=1}^4d_i^2
       -{cd_G\over S}u^2R^2-{c\over S}\eps_G\Gamma_*^2\\
 &-{c\over S}\{\sqrt{a_0u}\,U_*+\sqrt{a_0k}\,K_*\}^2.
\end{align*}
Define $G$ by
$g^{\mathsf T}Gg=g^{\mathsf T}D_*g+\mathcal E$.
This is the actual main Gram after retaining the positive excess
of the $N_x,J_C,J_\ell$ squares, splitting $K_A+c\ge1$ equally
between $A_\theta,T_\theta$, and retaining their nonnegative
difference. The averaged circuit identity and angular identity
used in this step are precisely those of the printed all-$H$
proof at
\href{https://github.com/anubhav70747/q222-finite-sdp/blob/0c1ce2e15ed067a6845464e2468559cea5274f2c/latex/certificates/c1_short.tex#L7550-L7740}
{lines 7550--7740}. This definition therefore prints every main-Gram
entry through one quadratic polynomial instead of twelve rows of
expanded rational functions.

\section{Rational coefficient bounds and fixed radicals}

Put $v=\sqrt u$. Every entry of $G-D_*$ is a polynomial in
$(v,w)$, with source coefficients, supported on
\[
 \{1,w,w^2,v,v^2,v^2w,v^4,v^4w\}
 \subset\{v^\alpha w^\beta:0\le\alpha\le4,0\le\beta\le2\}.
\]
The following elementary bounds provide a reproducible coefficient
estimate:
\[
 |R_x|,|A_y|\le32,\quad |\chi|\le64,\quad
 |B_x|\le48\eps,\quad |k_1|\le16.
\]
The coefficients of $E_{yy}$ on $(u^2,w^2)$ are bounded by $(4,4)$;
those of $E_{RR}$ on $(1,u^2,w^2)$ by $(64\sigma,16,4)$;
those of $E_{yR}$ by $(32\sigma,12,8)$.
These follow term by term from \eqref{eq:basic-bounds}--\eqref{eq:p-bounds};
in the constant terms use $C(r-1)$, never $C$ separately.

Here is a single positive-coefficient majorant which verifies
all remaining estimates. In it $s=\sqrt\eps$, all twelve root
letters are formal nonnegative indeterminates, and
$n=r'=J_C+J_\ell$. Define
\[
 f_1=u\Gamma_*+sd_1+ud_2+(u+s)d_3+sT_\theta,
\]
\[
\begin{split}
 f_2={}&512s^2w\Gamma_*+512s^4r'
  +512s^2(w+s)(d_2+d_3)\\
 &+(2048s^2+1024sw)T_\theta+16uwf_1.
\end{split}
\]
The absolute coefficient polynomial of $\mathcal E$ is bounded by
\begin{align*}
 \mathcal Q={}&(4u^2+4w^2)n^2
 +(64s^2+16u^2+4w^2)r'^2
 +(32s^2+12u^2+8w^2)nr'\\
 &+12N_xf_1+J_Cf_2+16uwN_xn+32uwN_xr'
      +48s^2wN_xA_\theta+32unA_\theta+64uA_\theta r'\\
 &+\tfrac52s^2r'^2+2vU_*r'
      +6(16s^2+u+s)\sum_i d_i^2+6u^2r'^2
      +(144u+48s)\Gamma_*^2\\
 &+48uU_*^2+96vU_*K_*+48sK_*^2.
\end{align*}
Indeed \eqref{eq:S-bounds} bounds $\sqrt{d_G}/S$, $\sqrt\rho/S$,
$d_G/S$, $\rho/S$ by one, while \eqref{eq:c-positive} gives
$\eps_C\le16\sigma+u+k$ and $\eps_G\le24u+8k$ when needed.
The two small coefficients in $F_2$ use
$|R_0q_0|\le2048\eps$ and \eqref{eq:R0-bounds}.

Expansion, or simply evaluation at one of this positive-coefficient
polynomial, gives coefficient $\ell^1$ norm $8590$. Its constant
testing part has norm $5030$ and every term of that part is
divisible by $s$. Consequently every source coefficient of every
nonconstant testing monomial in $G$ has absolute value below
$2^{20}$, and the source matrix $G(0,0)$ satisfies
\begin{equation}\label{eq:matrixbounds}
 |G_{ij}(0,0)-(D_*)_{ij}|\le2^{20}\sqrt\eps=2^{-60}.
\end{equation}
These estimates hold before any ordering of the tensor product is
assumed. They bound pure-source coefficients in their own field.

For $r=12$ the following exact quadratic-form decomposition uses
only fixed directions:
\[
\begin{split}
 G={}&\sum_i\left(G_{ii}-{11\over48}\right)e_ie_i^{\mathsf T}\\
 &+\sum_{i<j}{1/48+G_{ij}\over2}(e_i+e_j)(e_i+e_j)^{\mathsf T}\\
 &+\sum_{i<j}{1/48-G_{ij}\over2}(e_i-e_j)(e_i-e_j)^{\mathsf T}.
\end{split}
\]
Write each of its $144$ scalar channels as
\[
 q(c,t)=q^{(0)}(c)+\sum_{j=1}^{14}a_j(c)b_j(t),
\]
where $b_j$ run through the nonconstant monomials in the
$5\times3$ rectangle above. By \eqref{eq:matrixbounds}, for every
channel
\[
 q^{(0)}\ge c_0=1/256,\qquad |a_j|\le A_*=2^{30}.
\]
In fact the off-diagonal channels have constant term at least
$1/96-2^{-61}$ and the diagonal ones at least $13/48-2^{-60}$.
On the independent testing rectangle,
$0\le b_j\le\sqrt\delta=2^{-50}<\gamma$.
Hence the exact separated identity
\[
 q=(q^{(0)}-c_0)+\frac12\sum_{j=1}^{14}
 \bigl[(A_*+a_j)(\gamma+b_j)+(A_*-a_j)(\gamma-b_j)\bigr]
\]
has nonnegative factors in their own labelled fields. It proves
the required tensor SOS without an appeal to entrywise positivity
in a tensor ring or to a point-dependent real congruence.

All twelve roots are source-linear combinations of the fixed
testing functions
\[
 \mathcal V=\{X,Y,Z,W,uZ,wZ,uW,uwW,u^2W,
                u^{3/2}W,u^{1/2}W,
                  \operatorname{sgn}(W)\sqrt{uw\max(WZ,0)}\}.
\]
The last function includes both state-sign cases and is zero
when $WZ\le0$. Changing all state signs changes all twelve roots'
signs and preserves the quadratic form.
Therefore a fixed main-Gram dictionary is
\[
 \mathcal V\ \cup\ \{f\sqrt{\gamma+b_j},f\sqrt{\gamma-b_j}:
                         f\in\mathcal V,\ 1\le j\le14\}.
\]
It has $348$ displayed functions. All radicals are testing-only;
all remaining coefficients are source-only square roots.

Only seven nonconstant monomials actually occur. Dropping the zero
channels and replacing $\gamma$ by $1/(7\,2^{38})$ gives the
smaller displayed list of $12(1+14)=180$ functions. This optional
reduction changes no source estimate. The accompanying exact
polynomial check records the support as well as both norms.

The separate multiplier of the retained circuit remainder is
$\tfrac12-cd_G\eps_C/S$. Its constant term is at least
$1/2-6(16\eps+\sqrt\eps)>c_0$ and its only testing coefficient
is bounded by six. More explicitly, put
\[
 a_C=-cd_G/S,\qquad
 q_C^{(0)}=\tfrac12-(cd_G/S)(2/\lambda+k).
\]
Then its fixed-radical factorization is
\[
 q_C^{(0)}+a_Cu=(q_C^{(0)}-c_0)+13A_*\gamma
 +\tfrac12\{(A_*+a_C)(\gamma+u)
                          +(A_*-a_C)(\gamma-u)\}.
\]
Every displayed pure factor is nonnegative on the rectangle.
Its circuit remainder roots, and the other retained squares,
still require their own testing features; they are not counted
in $348$ or $180$.

\section{The two zero faces are present, and $H$ is unrestricted}

No inverse of $D$, $\beta$, $b$, $Q$ or a state coordinate was used
in the coefficient estimates or Gram decomposition. At $D=0$
the root $\Gamma_*$ is zero and all formulas remain literal.
At $\beta=0$ the same formulas apply. Testing axes are included.
The denominators $S,M,P,J,\sigma,k$ used above are nonzero on
the stated source domain; their potentially dangerous combinations
were bounded explicitly rather than omitted on a boundary.

For explicit nonemptiness set $x=\sigma=t$, $z=1$, $b=1/2$,
$h=Ht$, with $0<t\le\eps$. The numerators $D_N=2JPD$ and
$\beta_N=2\theta P\beta$ are exactly
\[
\begin{split}
 D_N={}&(3-H)t^2+(3/4-6H)t^3+(H^2+H/2)t^4\\
       &+(-7H^2/2+H/4)t^5+(H^3-H^2/2)t^6,\\
 \beta_N={}&(H-3/2)t+4Ht^2-Ht^3+4H^2t^4
                                      +(-H^3+H^2/2)t^5.
\end{split}
\]
Uniformly for $1\le H\le4$, the remainders beyond the first
terms have absolute values at most $(687/4)t^3$ and $156t^2$.
Thus continuity gives a $\beta=0$ source with $H\in[1,2]$ and
$D>0$, and a $D=0$ source with $H\in[2,4]$ and $\beta>0$.
On the first face $p\le7/8$ and
$\lambda=8DH+2(H^2-p)t\ge t/4$.
On the second, $\beta\theta\ge t/16$, so
$\lambda\ge1/8-4t>0$. All strict native gaps hold.

Finally take $\sigma=t$, $x=t^2/16$, $z=1$, $b=1/2$, $h=3/4$.
Then $H=3/(4t)$ is unbounded along $t\downarrow0$ and
\[
 D_N={-3t^2\over262144}
 (6t^5-61t^4+608t^3-7872t^2+8192t-126976)>0,
\]
\[
 \beta_N={3(3t^2+64)(2t^3+5t^2+64)\over16384}>0.
\]
Also $K_A=O(\eps)$ on this path by \eqref{eq:KA-bounds}, so
$c>0$ is represented. These are singular limits inside the
evaluated rectangle, not merely a nonsingular compact slice.

\section{Exact checks and scope}

The adjacent script \texttt{fixed\_box\_c1\_checks.py} verifies with
rational polynomial arithmetic the $F_x$ norm $37$, $F_\sigma$,
the majorant norms $8590$ and $5030$, its testing support, all
displayed rational margin comparisons, and the numerator formulas
and sign brackets for the two faces. Its JSON output records each
check. The inequalities above also supply the analytical proof;
the script is not a numerical fit or a replacement for it.

This supplies the previously missing fixed constants and fixed
testing radicals for the positive-angular main Gram. It does not
compile the other retained square families, the $\lambda\le0$
branch, complementary testing charts, other C1 source charts or
the global common-cover selection. None of those cases is silently
included in the stated feature count.
\end{document}
